\chapter*{Sommario}
\addcontentsline{toc}{chapter}{Sommario}

La consultazione delle fonti di osservazione della Terra richiede di distinguere
eventi naturali, acquisizioni satellitari, attivazioni cartografiche e cataloghi
di dataset. Questa tesi presenta Copernicus Data Workspace, un prototipo
conversazionale che interpreta richieste geografiche e coordina il recupero
di informazioni mantenendone riconoscibili provenienza e limiti.

Il sistema adotta un workflow LangGraph con componenti decisionali basati su
modelli linguistici e trasformazioni deterministiche. Il percorso di Discovery
interroga un catalogo di metadati NebulaGraph alimentato da harvester Marine,
STAC e Drought. Il percorso operativo consulta GDACS, Copernicus Emergency
Management Service e il catalogo STAC remoto; MCP delimita le integrazioni
GDACS e CEMS. Memoria conversazionale, artefatti e telemetria sostengono
rispettivamente i turni successivi, la presentazione dei risultati e la diagnosi.
Conteggi e aggregazioni delle visualizzazioni derivano dai record recuperati,
senza attribuire al modello il ruolo di sorgente dei valori.

La valutazione distingue controlli software, benchmark del workflow e verifica
dell'interfaccia. Il protocollo descritto combina controlli deterministici e
LLM-as-a-Judge, rendendo espliciti denominatori, dipendenze tra casi e limiti
del valutatore. Sul catalogo a grafo il sistema supera circa due casi su tre,
con il limite principale nell'ordinamento dei risultati; sulle acquisizioni
satellitari supera tutti i casi del benchmark. Il contributo è progettuale e implementativo:
non viene dimostrata una superiorità quantitativa rispetto a un agente singolo
o a una pipeline più semplice. Restano necessari confronti controllati,
validazione umana dei giudizi e prove conversazionali nel browser.

\clearpage
